\section{Discussion}\label{sec:discussion}

At median $\rho = \MedianRho{}$ the two methods are not the wildly divergent pair some of
the disagreement literature implies, but a third of each top-five list still differs. The
practical consequence is that reading the intersection of the two is not a robustness
check --- it silently drops about a third of each explanation --- and if a decision turns
on which features appear in a top-five list, the choice of explainer is part of that
decision and belongs in the record.

On this evidence the default should be \shap{} wherever a tree backend exists: on these
\NumCells{} cells it wins on every axis we measured except model-agnosticism. That
remaining axis is precisely \lime{}'s case, and it is not a small one.

The stability result is the one we would most want carried forward. Tree \shap{}'s re-run
variance is exactly zero, so a study measuring stability that way alone would report it as
perfectly stable; its cross-refit agreement was \RetrainShap{}, meaning the global ranking
moves when the model is refitted on a resample of the same training data. The instability
lives in the model, not the explainer, and no amount of determinism in the attribution
method removes it. Anyone reporting explanation stability should say which of the three
things they measured.

Freezing the protocol repaid itself three times. The additivity check caught a
background-construction error that violated \shap{}'s own additivity by one per cent on
LightGBM cells while every other cell looked healthy; a test requiring each interval to
contain its own point estimate caught a misapplication of BCa bootstrap to a non-smooth
statistic; and a test on the \rbo{} floor caught that raw \rbo{} is not comparable across
datasets of different width. None would have been visible in the results, and all three
are recorded in the protocol's amendment log.
